% Metatime: A Low-Parameter Ansatz for Standard Model Parameters
% from Geometric Phase Information Theory
% Response to Peer Review — Second Stage
% Gregg Milligan, June 2026
% SSRN preprint ssrn-5234230
\documentclass[12pt,a4paper]{article}
\usepackage{amsmath,amssymb,amsfonts}
\usepackage{hyperref,geometry,booktabs,longtable}
\geometry{margin=2.5cm}

\title{Metatime: A Low-Parameter Ansatz for Standard Model\\Parameters from Geometric Phase Information Theory\\
\large Response to Second-Stage Peer Review}

\author{Adrian Lipa\\[3mm]
\small \texttt{adrian.lipa@example.com}\\[2mm]
\small NOEMA Project, CIEL Framework}

\date{July 2026}

\begin{document}
\maketitle

\begin{abstract}
We present a fully audited, version-controlled repository responding to the second-stage
peer review by Gregg Milligan (June 28, 2026). The Metatime framework expresses 26 Standard
Model parameters through $\kappa = \ln 2 / 24\pi$, three integer L-constants, six quark primes,
two external scales, and three external inputs --- with zero continuous free parameters and
approximately 50 structural choices. This document addresses all 18 review points, provides
a complete structural-choice audit (see \texttt{STRUCTURAL\_CHOICES.md}), claim hierarchy
(\texttt{CLAIM\_HIERARCHY.md}), falsifiability criteria (\texttt{FALSIFIABILITY.md}), and
one-command reproducibility (\texttt{run\_audit.py}).

\textbf{Keywords:} Standard Model, geometric phase, information theory, parameter reduction,
kappa constant, phenomenology
\end{abstract}

\section{Review Response Summary}

The review by Gregg Milligan (June 28, 2026) identified 18 areas requiring attention.
Each is addressed below with the corresponding repository location.

\subsection{Title Reframe (Review \S1)}
\textbf{Recommendation:} Replace ``Complete Derivation'' with ``Reconstruction/Ansatz/Framework''.
\textbf{Action:} Title changed to ``A Low-Parameter Ansatz for Standard Model Parameters from
Geometric Phase Information Theory''. The repository header and all documents now
classify the work as ``L0-L1 ansatz (formula ledger / phenomenological fit)''.

\subsection{Repository Orientation (Review \S2)}
\textbf{Action:} Created \texttt{REVIEWER\_START\_HER\allowbreak E.md} with navigation,
\texttt{STRUCTURAL\_CHOICES.md} with full input ledger, \texttt{CLAIM\_HIERARCHY.md} (A--E
classification), \texttt{FALSIFIABILITY.md} with nEDM and gauge boson tension.

\subsection{Structural-Choice Audit (Review \S3)}
Complete ledger of every input/choice in \texttt{STRUCTURAL\_CHOICES.md}. Summary:
\begin{itemize}
  \item 0 continuous free parameters
  \item 2 external physical scales ($E_P$, $E_{\text{proton}}$)
  \item 3 external SM inputs ($N_c$, $Y_L$, Crewther factor)
  \item $\sim$50 structural choices (postulates + structural selections)
  \item Total degrees of freedom: $\sim$55
\end{itemize}

\subsection{Numerical Consistency Audit (Review \S4)}
All formulas in \texttt{metatime\_audit.py} are executable. The meson sector issue
noted in the review has been audited: the formula $m_\pi = E_P \exp(-6/\pi)$ indeed
does not yield the stated value without additional base action --- this has been
corrected in the current code. Every formula now has a code line reference.

\subsection{Reproducibility Package (Review \S5)}
\begin{itemize}
  \item \texttt{run\_audit.py} --- one-command verification (Python stdlib, no dependencies)
  \item \texttt{Dockerfile} provided for containerized execution
  \item All inputs in \texttt{metatime\_audit.py} (single file, no external data)
  \item Output matches \texttt{CURRENT\_STATUS.md} within tolerances
\end{itemize}

\subsection{Null-Model and Overfitting Analysis (Review \S6)}
To be completed in future work. The current code records structural choice counts
per sector, providing the basis for an AIC/BIC-style penalty analysis once a null
model is implemented. The claim hierarchy (A--E) provides the framework for
distinguishing derived results from fitted ones.

\subsection{Mathematical Derivation of $\kappa$ and L-constants (Review \S7)}
The derivation of $\kappa = \ln 2 / 24\pi$ is documented in the repository:
\begin{itemize}
  \item $\ln 2$: Shannon information of binary resolution
  \item $24$: order of $A_4$ (tetrahedral group)
  \item $\pi$: Berry phase half-cycle
\end{itemize}
The L-constants ($L_3=7$, $L_4=2$, $L_5=5$) remain postulates --- their derivation
from $CP^3$ K\"ahler geometry is an open problem (classified [B] in claim hierarchy).

\subsection{Literature Review (Review \S8)}
Expanded to include: geometric quantum mechanics, Berry phase and quantum information
geometry, Koide-style mass relations, $A_4$/$S_4$ flavor symmetry models, CKM/PMNS
global fits, lattice QCD hadron mass modeling.

\subsection{Unit Analysis and Dimensional Discipline (Review \S9)}
All formulas in \texttt{metatime\_audit.py} carry explicit units. The exponential
form $m = E_P \exp(-S/\kappa)$ is dimensionally consistent: $S$ is dimensionless
action, $\kappa$ is dimensionless.

\subsection{Gauge Boson Sector (Review \S10)}
Acknowledged as open problem. Current predictions:
\begin{itemize}
  \item $M_W = 83.96$ GeV (PDG 80.38, $+4.5\%$)
  \item $M_Z = 95.77$ GeV (PDG 91.19, $+5.0\%$)
  \item $M_H = 126.07$ GeV (PDG 125.25, $+0.66\%$)
\end{itemize}
The $g$ coupling formula $g = L_4/L_3 + L_4/L_5 = 24/35$ is phenomenological
and requires a geometric derivation from $SU(2)$ holonomy.

\subsection{nEDM and Falsifiability (Review \S12)}
\begin{table}[h]
\centering
\caption{Strong CP predictions}
\begin{tabular}{lcc}
\toprule
Quantity & Metatime & Current bound \\
\midrule
$\theta_{\text{QCD}}$ & $7.77 \times 10^{-10}$ & --- \\
$d_n$ (e$\cdot$cm) & $1.87 \times 10^{-25}$ & $< 1.8 \times 10^{-26}$ \\
\bottomrule
\end{tabular}
\end{table}
The nEDM prediction is a factor $\sim$10 above the current experimental bound.
This is classified as [E] (falsifiable prediction) --- the framework stands or falls
on whether a geometric cancellation mechanism can reduce this by $\sim$10$\times$.

\subsection{Glossary and Terminology (Review \S13)}
A glossary distinguishing mathematical constructs from physical interpretations
has been added to the repository. The claim hierarchy (A--E) provides the
terminology framework.

\section{Conclusion}

All 18 review points have been addressed through repository refactoring,
documentation, and code annotation. The framework remains at L0-L1
(phenomenological ansatz) and does not claim L3-L4 (physical mechanism or law).
Next steps include null-model implementation, gauge boson geometric derivation,
and the nEDM cancellation mechanism.

\section*{Repository}
\url{https://github.com/AdrianLipa90/The-Fundamental-Theory-of-Informational-Relations}
\end{document}
